\section{Introduction}
% 引言。

GPU resource management policies critically influence application performance.
% GPU 资源管理策略对应用性能有至关重要的影响。
These policies include memory placement, migration, and GPU task scheduling.
% 这些策略包括内存放置、迁移和 GPU 任务调度。
Recent research shows their substantial impact on Unified Virtual Memory (UVM) placement~\cite{wang2024suv,lin2025forest,kim2025dream,park2025helm} and GPU scheduling strategies~\cite{wang2024gcaps,fan2025gpreempt,shen2025xsched}.
% 近期研究表明，它们对统一虚拟内存 (UVM) 的放置和 GPU 调度策略有显著影响。
For instance, optimal eviction and prefetch policies can improve performance by up to 73\% under memory oversubscription~\cite{park2025helm,kehne2019etc,ganguly2019interplay}.
% 例如，在内存超配下，最优的逐出与预取策略可将性能提升多达 73\%。

Unfortunately, finding an optimal GPU resource management policy is challenging.
% 遗憾的是，寻找最优的 GPU 资源管理策略颇具挑战。
The optimal memory and scheduling policies vary significantly across workloads, hardware, and deployment contexts~\cite{kwon2023efficient,agrawal2024taming,lin2024towards,wang2016gunrock,shoeybi2019megatron,rasley2020deepspeed,pan2024survey,cao2023gpu}, with no one policy dominating the others.
% 最优的内存与调度策略随负载、硬件和部署环境的不同而差异显著，不存在一个策略在所有场景下都占优。
For example, aggressive memory prefetching improves latency for LLM MoE inference by reducing memory stalls~\cite{park2025helm,ganguly2019interplay}, but increases the latency of graph analytics queries because it over-saturates memory bandwidth~\cite{ganguly2019interplay}.
% 例如，激进的内存预取通过减少内存停顿改善了 LLM MoE 推理的延迟，但会因内存带宽过饱和而增加图分析查询的延迟。
Additionally, throughput-oriented scheduling improves average latency of vector search but degrades tail latency for interactive inference~\cite{agrawal2024taming,fan2025gpreempt}.
% 此外，面向吞吐量的调度改善了向量搜索的平均延迟，却恶化了交互式推理的尾延迟。

%\arq{I don't love this sentence structure.}
We argue that the community should employ AI agents to automatically optimize programmable GPU resource management policies.
% 我们认为，社区应当利用 AI 智能体来自动优化可编程的 GPU 资源管理策略。
AI-driven exploration can handle the vast, workload-specific policy space that makes manual tuning challenging.
% AI 驱动的探索能够应对庞大且随负载而异的策略空间，而这正是手动调优的难点。
% Vast, workload-specific policy space make manual tuning costly but suitable for AI-driven exploration. 
Agents can now generate complex code and optimize policies in domains such as compiler heuristics~\cite{alphaevolve2025}, datacenter scheduling, LLM serving~\cite{hamadanian2025glia,cheng2025barbarians}, training strategies~\cite{ding2025asap}, and CPU scheduling~\cite{dwivedula2025policysmith,zheng2025schedcp}.
% 如今，智能体已能生成复杂代码，并在编译器启发式、数据中心调度、LLM 服务、训练策略和 CPU 调度等领域优化策略。
In these use-cases, agents are generating fully fledged policy programs rather than simply tuning configuration parameters.
% 在这些场景中，智能体生成的是完整的策略程序，而非简单地调节配置参数。

We derive three properties that systems should provide to enable AI agents to optimize programmable GPU resource management policies effectively without sacrificing system health.
% 我们提炼出系统应提供的三项特性，使 AI 智能体能够在不损害系统健康的前提下有效优化可编程的 GPU 资源管理策略。
Namely, a system should support programmable policy interfaces that are safe, dynamic, and full-stack.
% 具体而言，系统应支持安全、动态且全栈的可编程策略接口。
A policy interface is safe if it ensures that new policies cannot cause the system to crash; dynamic if it supports policy updates without requiring the system and applications to be restarted; and full-stack if it supports cross-application and device-internal visibility, as well as low-level hardware control.
% 所谓安全，是指保证新策略不会导致系统崩溃；所谓动态，是指支持策略更新而无需重启系统和应用；所谓全栈，是指既支持跨应用和设备内部的可见性，又支持底层硬件控制。

Existing programmable GPU resource management techniques fail to simultaneously meet these requirements.
% 现有的可编程 GPU 资源管理技术无法同时满足这些要求。
User-space frameworks~\cite{ng2023paella,shen2025xsched,gim2025pie,chen2025ktransformers} are safe and dynamic but are not full-stack since policies lack cross-application visibility and hardware control.
% 用户态框架安全且动态，但并非全栈，因为策略缺乏跨应用可见性和硬件控制能力。
Other approaches modify the OS kernel driver~\cite{kato2011timegraph,kato2012gdev,kernelintercept,fan2025gpreempt,wang2024gcaps}, which fails to meet all three of the requirements.
% 另一些方案修改 OS 内核驱动，也无法同时满足这三项要求。
%full-stack at the host level but neither safe nor dynamic, and still lacks GPU-internal observability.

In this paper, we propose \sys, a new system that provides an OS interface for GPU resource management policy that is safe, dynamic, and full-stack.
% 本文提出 \sys，一个提供安全、动态且全栈的 GPU 资源管理策略 OS 接口的新系统。
Similar to prior work on resource management~\cite{schedext-docs, zussman2025cache_ext, bachl2021flow}, \sys{} supports eBPF-based policies.
% 与以往的资源管理工作类似，\sys 支持基于 eBPF 的策略。
Users write callback functions that manage system resources and are called on user-defined events (e.g., when the system reaches a specific code location, or periodically).
% 用户编写管理系统资源的回调函数，它们由用户自定义的事件触发（例如系统到达特定代码位置时，或周期性地触发）。
\sys{}'s use of eBPF ensures safety, since the system uses an eBPF-style verifier before running new policies, and dynamism, since it can replace policies without service interruptions.
% \sys 对 eBPF 的使用保证了安全性——系统在运行新策略前会进行 eBPF 风格的验证——以及动态性——它可以在不中断服务的情况下替换策略。
To provide a full-stack interface, \sys{} supports eBPF callbacks that execute within the OS kernel driver and GPU device.
% 为提供全栈接口，\sys 支持在 OS 内核驱动和 GPU 设备内执行的 eBPF 回调。

\sys models GPU resource management as transitions on resource state machines (\S\ref{sec:design-interfaces}).
% \sys 将 GPU 资源管理建模为资源状态机上的状态转移。
State machines define valid states and transitions: GPU memory chunks transition between host memory and device memory through eviction and prefetching; compute contexts transition through states such as running, ready, and finished via scheduling and preemption operations.
% 状态机定义合法的状态与转移：GPU 内存块通过逐出和预取在主机内存与设备内存之间转移；计算上下文通过调度与抢占操作在运行、就绪、完成等状态之间转移。
eBPF policies specify the transitions that the system should take on its GPU resources.
% eBPF 策略指定系统应在其 GPU 资源上执行哪些转移。



%, with a SIMT-aware verifier and device JIT runtime. Bring these up in the implementation paragraph.

Alas, applying the CPU eBPF programming model to GPU resource management is not straightforward.
% 遗憾的是，将 CPU 上的 eBPF 编程模型直接套用于 GPU 资源管理并非易事。
The GPU is a discrete processor with its own execution model and resources, and policy decisions made in the host take effect on a separate device.
% GPU 是拥有自身执行模型和资源的独立处理器，主机上做出的策略决策要在另一台独立设备上生效。
In particular, we identify three challenges: first, cross-device operations like migrating data or dispatching commands take milliseconds, too slow for synchronous callbacks and too costly for purely reactive decisions.
% 具体而言，我们识别出三项挑战：第一，迁移数据或派发命令等跨设备操作需要毫秒级时间，对同步回调而言太慢，对纯粹被动响应式决策而言代价过高。
Second, unlike independent subsystems on the CPU, limited device memory and bandwidth further couple memory and scheduling, as decisions about data placement directly affect which workloads can execute.
% 第二，与 CPU 上相互独立的子系统不同，有限的设备内存和带宽进一步耦合了内存与调度，因为数据放置的决策直接决定哪些负载可以执行。
Third, GPU applications employ OS kernel bypass, which requires \sys{} to support eBPF callbacks from events in user-space applications.
% 第三，GPU 应用采用 OS 内核旁路，这要求 \sys 支持来自用户态应用事件的 eBPF 回调。
Traditional eBPF does not allow such callbacks to modify OS state required to implement resource management policies because of safety concerns.
% 出于安全考虑，传统 eBPF 不允许这类回调修改实现资源管理策略所需的 OS 状态。

\sys{} addresses these challenges through a novel execution model for eBPF-based resource management that separates the effect of a callback from its triggering event.
% \sys 通过一种面向基于 eBPF 的资源管理的新型执行模型来应对这些挑战，将回调的效果与其触发事件解耦。
Unlike prior eBPF-based systems (e.g., sched\_ext~\cite{schedext-docs}, cache\_ext~\cite{zussman2025cache_ext}, and XDP~\cite{bachl2021flow}), where each callback makes a single synchronous decision at a fixed event, \sys{} allows callbacks that asynchronously transition resources at user-defined events.
% 不同于以往的 eBPF 系统（如 sched\_ext、cache\_ext 和 XDP），它们的每个回调在固定事件处做出单一的同步决策，\sys 允许回调在用户自定义事件处异步地转移资源。
In particular, \sys callbacks invoke transitions when a given event occurs, and the system asynchronously transitions the resource between states.
% 具体而言，\sys 回调在给定事件发生时被触发，系统随后异步地在状态之间转移资源。
\sys{} callbacks can initiate multiple state transitions within a single event.
% \sys 回调可以在单个事件中发起多次状态转移。
Finally, \sys{} supports user-defined events that can execute policy code from anywhere in the GPU stack (application, OS kernel driver, or GPU device).
% 最后，\sys 支持用户自定义事件，可在 GPU 栈的任意位置（应用、OS 内核驱动或 GPU 设备）执行策略代码。
Critically, to enable safety with user-defined events, \sys{} policies only initiate predefined transitions; they cannot create invalid transitions that could corrupt resources.
% 关键在于，为了在使用用户自定义事件时仍保证安全，\sys 策略只能发起预定义的转移，无法创建可能破坏资源的非法转移。

% To ensure safety, policies only influence which transitions occur and when, but cannot create invalid transitions.


We implement \sys{} on Linux by extending the NVIDIA open GPU OS kernel modules~\cite{nvidia-open-gpu} and providing an eBPF SIMT verifier and JIT runtime for GPU devices.
% 我们在 Linux 上实现 \sys，扩展 NVIDIA 开源 GPU OS 内核模块，并为 GPU 设备提供 eBPF SIMT 验证器和 JIT 运行时。
We evaluate \sys{} across LLM inference, GNN training, vector search, and multi-tenant scenarios.
% 我们在 LLM 推理、GNN 训练、向量搜索和多租户场景下评估 \sys。
\sys-based policies improve throughput by up to $4.8\times$ over framework offloading and $1.76\times$ over application-tuned baselines, and reduce tail latency by up to $2\times$, with low overhead and no application modifications or restarts.
% 基于 \sys 的策略相比框架卸载最高提升 $4.8\times$ 吞吐量，相比应用调优的基线最高提升 $1.76\times$，并将尾延迟最多降低 $2\times$，同时开销很低，且无需修改或重启应用。
These gains come from the policies, not from an intrinsic BPF speedup: expressing the same policy through \sys{} costs 0.7--3.3\% in mechanism overhead, and \sys{} reimplementations of seven published resource-management systems (including MoE-Infinity, XSched, and GPREEMPT) match the original ad-hoc implementations within a few percent (\S\ref{sec:policy-mechanism}).
% 这些收益来自策略本身，而非 BPF 的固有加速：通过 \sys 表达同一策略的机制开销仅为 0.7--3.3\%，且 \sys 对七个已发表资源管理系统（包括 MoE-Infinity、XSched 和 GPREEMPT）的重实现与原有的临时实现相差仅在几个百分点以内（\S\ref{sec:policy-mechanism}）。
We show that \sys{} enforces safety by using AI agents to generate all of the 59 evaluated eBPF policies, including eviction ordering, prefetch strategies, scheduling logic, and device-side work-stealing.
% 我们通过让 AI 智能体生成全部 59 个受测 eBPF 策略（包括逐出顺序、预取策略、调度逻辑和设备端工作窃取）来证明 \sys 的安全性。
\sys{} ensured that there were zero OS kernel panics throughout this process.
% 在整个过程中，\sys 保证了零次 OS 内核恐慌。
In summary, our contributions are:
% 总而言之，我们的贡献如下：

\begin{itemize}[leftmargin=*,itemsep=0pt]
  \item an execution model for eBPF-based resource management that separates the effect of a callback from its triggering event.
  % 一项面向基于 eBPF 的资源管理的执行模型，将回调的效果与其触发事件解耦。
  %We design a GPU-specific OS policy interface that models GPU resources as state machines whose asynchronous state transitions are directed by verified eBPF programs across the host-device boundary.
  \item \sys, a system that provides a programmable GPU resource management policy interface that is safe, dynamic, and full-stack.
  % \sys：一个提供安全、动态且全栈的可编程 GPU 资源管理策略接口的系统。
  %We implement \sys, an eBPF runtime spanning host driver and GPU device, with a SIMT-aware verifier for device-side eBPF execution.
  \item an evaluation of \sys across four GPU workloads, achieving up to $4.8\times$ throughput improvement and $2\times$ tail latency reduction, with all 59 policies generated by AI agents and no OS kernel panics; matched reimplementations of seven published policies isolate the mechanism's cost from the policies' benefits.
  % 在四类 GPU 负载上对 \sys 的评估：吞吐量最高提升 $4.8\times$，尾延迟最多降低 $2\times$，全部 59 个策略由 AI 智能体生成，且无 OS 内核恐慌；对七个已发表策略的匹配重实现将机制代价与策略收益分离。
\end{itemize}
